Zero-shot text-to-speech (TTS) can clone a voice from seconds of audio. 
Yet the quality comes from billion-parameter diffusion transformers whose size and multi-step sampling put them out of reach of the on-device, assistive, and real-time applications that need them most.
Post-training quantization (PTQ) shrinks such models without retraining, but it has been built for language and vision, not speech. 
At 4 bits, quantization harms intelligibility far more than timbre, so similarity scores stay high while the model becomes harder to understand. 
We ask why quantization hurts speech, where the damage lives, and whether it can be repaired without retraining. 
To find out, we run the first 4-bit weight-and-activation (W4A4) PTQ study of a TTS diffusion transformer at 1B and 3.5B. 
Intelligibility damage traces to the weight quantization, timbre damage to the activation quantization, and the precision-allocation rules from vision reverse in speech. 
Because timbre damage lives in the activations, we keep only the late-step feed-forward activations at full precision, one sixth of the budget. 
The method recovers timbre best at both scales, at no cost to intelligibility. 
This brings high-quality TTS within reach of the real-time and on-device applications that need it.
